% 5.4 — Le portail de décision
\section{Le portail de décision}
\label{sec:portail}

\subsection{Principe}
Le portail évalue un corpus de \nbcorpus{} commandes avec le code avant et après une
modification, et compare pour chaque
ligne la décision, la règle, la cible et la capture. Il ne sait pas quelle décision est juste ;
il voit chaque décision qui \emph{change}. Une ligne qui bouge doit être expliquée : erreur
corrigée, arbitrage appliqué, ou simple changement d'étiquette. La référence est ensuite
réenregistrée.

\subsection{Les mouvements, partie par partie}
Le tableau~\ref{tab:portail} reprend ce que chaque message de commit déclare. Les parties qui
ne devaient rien changer --- le moteur en ombre, les canaux ask/deny, plusieurs familles --- ne
changent rien ; les autres changent des lignes \emph{nommées et revues}.

\begin{table}[htbp]
\centering
\caption{Mouvements du portail de décision déclarés par partie (messages de commit)}
\label{tab:portail}
\small
\begin{tabularx}{\linewidth}{@{}L{2.6cm}L{2.3cm}Y@{}}
\toprule
\textbf{Partie} & \textbf{Lignes} & \textbf{Nature} \\
\midrule
B, C1, C2a, SYS1 & 0 & moteur en ombre, filets reliés : aucune décision ne change \\
S2+M1 & 124 & 8 erreurs corrigées, 18 nouvelles captures, 86 nouvelles autorisations (rien
  d'existant n'était détruit) --- le moteur prend la main \\
M2.0+M2.1 & 28 & suppression SQLite désormais restaurée ; chemins absents hors projet autorisés \\
T & 0 & sur les deux portails (Linux et Windows) \\
M2.2a (git) & 26 & lignes classées, référence réenregistrée \\
M2.2b1 (infra) & 4 & 3 corrections, 1 arbitrage (R1) \\
M2.2b2a/b (cloud) & 1 / 0 & une étiquette seulement \\
M2.2b3 & 37 & comptage des surfaces ; \cmd{compose down} et \cmd{docker push} escaladent \\
ASK & 0 & le canal s'ajoute à l'escalade sans changer la décision \\
M2.2b4a (envois) & 24 & 16 comptages, 3 étiquettes \cmd{ssh}, 5 décisions arbitrées \\
M2.2b4b1/b2 & 12 / 9 & PaaS (\cmd{firebase deploy} escalade), SaaS et migrations \\
M2.2c, d, e & 0 & système, effet LATER, préfiltre généré \\
M2.2f & 10 & derniers outils signalés \\
M2.3a & 6 & toutes plus strictes \\
M2.3b & 65 & plus strictes ou arbitrées (Rb1--Rb4) ; référence Windows enregistrée à part \\
C3 & 57 & 30 étiquettes, 2 plus strictes, 25 plus permissives : 3 arbitrées (PS7), 4 faits
  PowerShell vérifiés, 18 dont la cible n'existait pas (§\ref{sec:portail-aveugle}) \\
Répertoire d'essai & 29 + 36 & l'instrument, pas l'outil : les fichiers manquants ajoutés en
  deux passes (§\ref{sec:portail-aveugle}) \\
C3b & 4 / 12 & Linux / Windows : 2 captures gagnées, 1 plus stricte ; sous Windows, 9
  redirections \cmd{/tmp} intraduisibles \\
\bottomrule
\end{tabularx}
\end{table}

\begin{pourquoi}[Pourquoi réenregistrer la référence ?]
Après revue, les mouvements deviennent la nouvelle norme : la partie suivante est comparée à
\emph{elle}, et non à un état ancien. Chaque ligne figée dans \fichier{tests/fixtures/} garde la
trace de l'arbitrage qui l'a fait bouger, ce qui permet de retrouver, des semaines plus tard,
pourquoi une commande escalade.
\end{pourquoi}

\subsection{Quand le portail ne suffit plus : le corpus élargi (M2.3b)}
M2.3b retirait à l'ancien chemin sa dernière ligne POSIX : un changement qui touche toutes les
commandes non lues. Le portail de \nbcorpus{} commandes ne suffisait plus à le valider. Un
\textbf{corpus élargi de \nbcorpuslarge{} commandes} --- le portail, toutes les lignes figées des
\textit{fixtures} et tous les exemples de la table de règles --- a été rejoué avant et après. Il a
révélé des régressions que le portail ne voyait pas, comme \cmd{kubectl exec \ldots{} rm -rf} ou
\cmd{npx prisma migrate reset}, corrigées avant la publication.

Chaque ligne devenue plus permissive y est soit une \emph{décision arbitrée}, soit un \emph{fait
vérifié} : la suppression d'un fichier absent (rien n'est détruit), un texte « rm -rf » dans les
métadonnées d'un envoi (une donnée, pas du code), un outil de migration qui crée seulement un
fichier local.

\subsection{L'instrument aveugle}
\label{sec:portail-aveugle}
En préparant C3b, l'examen des 25 lignes devenues plus permissives avec C3 a révélé que 18
d'entre elles ne disaient rien du moteur : elles visaient des fichiers que le répertoire d'essai
du portail ne créait jamais (\fichier{test.txt}, \fichier{victim.txt}, puis \fichier{notes.txt}
et \fichier{x}). Une suppression d'un fichier absent ne détruit rien ; ces lignes passaient donc
quelle que soit la décision du moteur. Sur ces lignes, le portail ne pouvait pas échouer.

La correction a porté sur l'instrument, pas sur l'outil : les fichiers manquants ont été ajoutés
en deux passes, et les deux références ont été réenregistrées. 29 puis 36 lignes ont bougé sur
chaque système, et \emph{aucune} ligne ne nommant pas l'un de ces fichiers n'a bougé. Ces lignes
capturent désormais un vrai fichier : une régression sur elles ferait enfin échouer le portail.
Deux lignes restent sans portée (\cmd{T=\$(cat x); rm \$T}, car \fichier{x} contient un texte
qui ne nomme aucun fichier), et quatre cibles du corpus restent absentes du répertoire d'essai
(\fichier{dist}, \fichier{proj}, \fichier{canary.txt}, \fichier{important.txt}) : les lignes qui
les visent ne peuvent toujours pas échouer. Le point est connu et assumé.

\begin{pourquoi}[Pourquoi c'est un résultat de méthode]
Un test vert ne prouve rien tant qu'on n'a pas montré qu'il \emph{peut} échouer. C'est
l'exigence du test qui échoue sur le code d'avant (§\ref{sec:fail-old}), appliquée cette fois à
l'instrument de mesure lui-même. Le défaut n'était visible ni dans le code ni dans les
résultats : seulement en se demandant pourquoi une ligne bougeait.
\end{pourquoi}

\subsection{Deux références, une par système}
Le portail a été réenregistré séparément pour Windows : 21 lignes ne bougent que sous Windows ---
des chemins POSIX absolus impossibles à traduire en chemins Windows, donc non compris ---, et
toutes y sont plus strictes ou de même décision. C3b a répété la leçon : 4 lignes bougent sous
Linux, 12 sous Windows. Les 9 lignes propres à Windows ont la même forme,
\cmd{powershell.exe -Command "Remove-Item -Force victim.txt" > /tmp/log} : le moteur lit la
suppression, mais la redirection vers \fichier{/tmp} ne peut pas être traduite en chemin
Windows ; la ligne n'est donc pas entièrement lue et escalade (M3), alors que la même ligne avec
\cmd{> out.log} reste réversible. Avant C3b, les anciennes règles ignoraient la redirection et
capturaient. Sous Linux, ces lignes sont identiques : validée sur un seul système, la partie
serait passée sans que personne ne le voie. Au commit de référence, le portail ne montre
aucune ligne différente sous Windows, et une seule sous Linux, qui dépend des fichiers présents
sur la machine.
